\chapter{Projekt bazy danych}
\label{cha:baza}

\section{Opis modelu danych}
\label{sec:modelDanych}

Wszystkie informacje o użytkownikach, grupach, testach, pytaniach, podejściach i odpowiedziach są
przechowywane w relacyjnej bazie danych PostgreSQL. Najważniejsze tabele przedstawiono
w tabeli~\ref{tab:tabele}. Można je podzielić na trzy obszary: konta użytkowników i grupy, testy
wraz z pytaniami oraz podejścia uczniów wraz z ich odpowiedziami.

\begin{table}[htbp]
  \centering
  \caption{Najważniejsze tabele bazy danych aplikacji.}
  \label{tab:tabele}
  \begin{tabularx}{\linewidth}{>{\bfseries}l X}
    \toprule
    Tabela & Przeznaczenie \\
    \midrule
    Users & Konta użytkowników (wspólna tabela nauczycieli i uczniów, rozróżnianych
            dyskryminatorem); login, skrót hasła, imię i nazwisko. \\
    Groups & Grupy uczniów (nazwa, opis). \\
    StudentGroups & Tabela łącząca uczniów z grupami (relacja wiele-do-wielu). \\
    Quizzes & Testy (nazwa, przedmiot, limit czasu, próg, liczba podejść, widoczność). \\
    Questions & Pytania (wspólna tabela pytań zamkniętych i otwartych); treść, punkty, kolejność. \\
    AnswerOptions & Warianty odpowiedzi pytań zamkniętych; treść, flaga poprawności. \\
    TestGroups & Tabela łącząca testy z grupami (relacja wiele-do-wielu). \\
    QuizAttempts & Podejścia uczniów do testów; wynik, procent, status, czas, data. \\
    StudentAnswers & Odpowiedzi ucznia (wspólna tabela odpowiedzi zamkniętych i otwartych). \\
    SelectedOptions & Warianty zaznaczone przez ucznia w odpowiedzi zamkniętej. \\
    \bottomrule
  \end{tabularx}
\end{table}

\section{Diagram ERD}
\label{sec:erd}

Centralnymi tabelami modelu są \texttt{Quizzes} (testy) oraz \texttt{Users} (użytkownicy). Test
posiada pytania i podejścia, a podejście ucznia gromadzi jego odpowiedzi. Pełny schemat powiązań
przedstawia diagram ERD na rysunku~\ref{rys:erd}.

\figscreen{erd.png}{0.95}{Diagram ERD bazy danych aplikacji.}{rys:erd}

\section{Opis relacji i kluczy}
\label{sec:relacje}

Każda tabela posiada klucz główny typu całkowitego (\texttt{IDENTITY}). Główne relacje:

\begin{itemize}
  \item \texttt{Quizzes} - \texttt{Questions} oraz \texttt{Quizzes} - \texttt{QuizAttempts}: relacje
        jeden-do-wielu (test ma wiele pytań i wiele podejść).
  \item \texttt{ClosedQuestion} - \texttt{AnswerOptions}: pytanie zamknięte ma wiele wariantów
        odpowiedzi.
  \item \texttt{QuizAttempts} - \texttt{StudentAnswers} - \texttt{SelectedOptions}: podejście ma wiele
        odpowiedzi, a odpowiedź zamknięta - wiele zaznaczonych wariantów.
  \item \texttt{StudentGroups} łączy uczniów i grupy w relacji wiele-do-wielu (klucz złożony
        \texttt{StudentId} + \texttt{GroupId}); analogicznie \texttt{TestGroups} łączy testy z grupami.
\end{itemize}

\section{Normalizacja}
\label{sec:normalizacja}

Baza danych spełnia trzecią postać normalną (3NF) - każda tabela przechowuje dane dotyczące jednego
pojęcia, wartości w kolumnach są atomowe, a powtarzające się informacje zostały zastąpione kluczami
obcymi. Relacje wiele-do-wielu (uczeń-grupa oraz test-grupa) rozwiązano za pomocą tabel łączących
(\texttt{StudentGroups}, \texttt{TestGroups}), co eliminuje powielanie danych. Dziedziczenie klas
odwzorowano strategią Table-per-Hierarchy, w której cała hierarchia mieści się w jednej tabeli,
a typ rekordu wskazuje kolumna dyskryminatora.

\section{Migracje bazy danych}
\label{sec:migracje}

Zmiany schematu bazy są zapisane w postaci migracji Entity Framework Core w katalogu
\texttt{Migrations}. Przy starcie aplikacja wywołuje metodę \texttt{Database.Migrate()}, która tworzy
bazę i nakłada wszystkie migracje. Kolejne migracje wprowadzały m.in. obsługę liczby podejść, zmianę
typu daty oddania na pełny znacznik czasu, unikalny indeks loginu oraz haszowanie haseł istniejących
kont. Dzięki migracjom schemat jest wersjonowany, a baza odtwarzalna bez ręcznego pisania skryptów SQL.

\section{Mechanizmy zapewniające integralność danych}
\label{sec:integralnosc}

\begin{itemize}
  \item \textbf{Integralność referencyjna} - klucze obce na wszystkich powiązaniach. Tam, gdzie dane
        podrzędne nie mają sensu bez nadrzędnych, zastosowano usuwanie kaskadowe (usunięcie testu
        usuwa jego pytania i podejścia); tam, gdzie usunięcie mogłoby naruszyć już oddane odpowiedzi,
        zastosowano regułę ograniczającą, a operacje usuwania wykonywane są we właściwej kolejności
        (najpierw rekordy zależne).
  \item \textbf{Unikalność} - unikalny indeks na loginie użytkownika gwarantuje brak dwóch kont o tej
        samej nazwie.
  \item \textbf{Spójność dziedzinowa} - walidacja po stronie aplikacji nie dopuszcza do zapisu testu
        niekompletnego (np. pytania bez poprawnej odpowiedzi).
\end{itemize}
